% ----------------------------------------------

\chapter{Introduction}

La phase de décollage constitue l'une des phases les plus importantes du vol d'un aéronef, concentrant des contraintes physiques, aérodynamiques majeures. \\

L'optimisation des performances aéronautiques repose sur la détermination précise de la distance de roulement au sol ($s_{\text{LO}}$) nécessaire au décollage d'un appareil. Cette distance dépend d'une interaction  entre les performances du système propulsif, les caractéristiques aérodynamiques de l'aéronef , les conditions de surface de la piste et l'état thermodynamique local de l'atmosphère (pression, température, densité). \\

De ce fait l'objectif de ce projet est de modéliser, de simuler et d'analyser la cinématique d'un avion lors de sa course au sol. Pour comprendre les forces qui lui sont appliquées, notamment la traînée et la portance qui évoluent de façon quadratique avec la vitesse. \\

Pour mener à bien cette étude, nous nous appuierons sur des modèles mathématiques  issus de la mécanique du vol, traduits sous forme de codes de calcul en environnement Python. Cette approche permet de confronter plusieurs méthodologies de résolution afin d'analyser finement le comportement aérodynamique et propulsif de divers aéronefs représentatifs de l'aviation aujourd'hui.


\vspace{1cm}
%============================================================================================================
%
%================================================================================================================



\section{Compétence à acquérir}
L'analyse de ce rapport permettra de développer les compétences suivantes :
\begin{itemize}[itemsep=0pt]
	\item \textbf{Comprendre :} la disposition des forces (comment elle affecte notre avion ) lors d'une montée stable et l'importance de l'angle de montée ($\gamma$)
	\item \textbf{Calculer  :} la vitesse ascensionnelle(Vc) en utilisant la notion de puissance excédentaire
	\item \textbf{l'influence :} de la densité de l'air et l'altitude sur les performances des turboréacteur (aussi des avions à hélices )

\end{itemize}


\section{ Outils utilisés  pour ce rapport}

	Pour mener cette étude, nous avons utilisé
	\begin{itemize}
		\item  Les modèles mathématiques de performance de montée proviennent du cours de mécanique du vol, issu de différents livres.
		\item Nous utilisons des applications comme VS Code pour la programmation python des codes afin de réaliser des analyses graphiques.
	\end{itemize}

	

	\vspace{1cm}
	
	
	
	%============================================================================================================
	%PARAMETRE INFLUENCANT LA DISTANCE DE DECOLLAGE 
	%================================================================================================================
	
	
	
	
	
	\section{ Les éléments qui influencent la distance de décollage }
	
	Lors d'un décollage on remarque quatre facteurs qui influencent la capacité de décollage des avions : les conditions atmosphériques, le poids de l'avion, les vents dominants et l'état de la piste. Ces éléments impactent la distance de décollage et la performance de l'appareil, rendant essentiel pour les pilotes de comprendre leur influence pour assurer la sécurité.
	
	
	
	\begin{itemize}
		\item \textbf{Les conditions atmosphériques}, telles que la température ambiante, affectent considérablement les performances de décollage des aéronefs en réduisant la puissance du moteur.
		
		
		Comme altitude-densité (pression atmosphérique) une pression plus basse (haute altitude ou basse pression barométrique) réduit la densité de l'air. \\
		
		Conséquence : Cela diminue à la fois la portance de l'aile et la poussée du moteur (en raison de la diminution de l'oxygène pour la combustion). \\
		
	
		
		
		\item\textbf{Le poids} de l'avion augmente la distance de décollage en raison de l'inertie plus importante, nécessitant ainsi une plus grande longueur de piste pour atteindre la hauteur de sécurité. Le poids affecte également l'angle de montée, ce qui nécessite une distance horizontale plus importante pour atteindre la hauteur de franchissement d'obstacle. \\
		
		Le poids total compte, mais l'endroit où il est placé aussi.
		Un avion centré trop "avant" (poids sur le nez) demandera une force beaucoup plus grande sur la gouverne de profondeur pour lever le nez.\\
		
		Conséquence: Cela augmente la vitesse de rotation ($V_r$) nécessaire et donc la distance de roulement diminue. À l'inverse, un centrage trop arrière rend l'avion instable.
		
		\item\textbf{Les vents dominants} jouent un rôle essentiel, car un vent de face peut réduire la distance de décollage, tandis qu'un vent arrière l'augmente, affectant la vitesse au sol.
		
		
		\begin{equation}
			s_{vent} \approx s_{sans vent} \left( \frac{V_{LOF} - V_{vent}}{V_{LOF}} \right)^2
		\end{equation}
		
		\textbf{Les conditions de la piste} et sa pente sont également significatives, car une piste parfaitement lisse n'existe pas, ce qui influence la friction et la distance de décollage.
		La friction entre les pneus de l'avion et la surface de la piste augmente avec l'eau, la neige ou la boue, ce qui complique le décollage.
		La pente de la piste influence aussi  l'accélération, une piste en montée augmente la distance de décollage, tandis qu'une pente descendante la réduit.\\
		
		
		\end{itemize}
			
			
			\subsection{Autres facteurs }
			Mais lors du décollage, d'autres facteurs sont cruciaux :\\
			
			- La configuration optimale des volets, déterminée lors des essais de l'appareil a un impact direct sur les performances de montée.\\
			
			- L'effet de sol: Ce phénomène aérodynamique se produit lorsque l'avion se trouve en dessous de son envergure. La proximité du sol réduit les tourbillons marginaux, ce qui diminue considérablement la traînée induite. Cela peut provoquer un décollage prématuré. Cependant, une fois sorti de l'effet de sol (quelques dizaines de pieds), la traînée augmente soudainement et l'avion peut s'enfoncer si la vitesse est insuffisante.\\
			
			\begin{equation}
				D = \frac{1}{2}\rho V^2 S \left( C_{D_0} + \phi \frac{C_L^2}{\pi AR e} \right)
			\end{equation}
			
			\begin{itemize}
				\item Au niveau de la piste, $\phi$ se situe entre $0{,}5$ et $0{,}7$ (forte réduction de traînée).
				\item Lorsque l'avion s'élève à plus d'une envergure du sol ($h > b$), $\phi \to 1$ (retour à la traînée normale).
			\end{itemize}
			
			
			\textbf {L'utilisation des servitudes (Prélèvement d'air) } \\
			
			C'est un facteur  qui est souvent oublié mais qui agit sur la performance moteur.
			Comme l'utilisation du dégivrage (anti-ice) ou de la climatisation prélève de l'air comprimé directement sur les moteurs.
			De ce fait, cela réduit la poussée maximale disponible pour l'accélération. En conditions de performance limite, les pilotes coupent souvent la climatisation pour le décollage.\\
			
			
				
			\textbf{La vitesse de rotation ($V_r$) } \\
			
			La manière dont le pilote manipule l'avion influence la distance réelle.\\
			En effet : Si le pilote tire sur la manche trop tôt (rotation précoce), il augmente la traînée avant d'avoir la vitesse nécessaire, ce qui prolonge la distance au sol.
			La vitesse de sécurité au décollage ($V_2$) : C'est la vitesse que l'avion doit atteindre à la "hauteur de franchissement d'obstacle" (screen height) pour garantir que, même en cas de panne moteur, la trajectoire reste sûre.
			
			
		
		
			\subsection*{Résumé :}
			
			\begin{itemize}
				\item \textbf{Conditions Atmo / Servitudes:} Font baisser $T$ et $\rho$.
				\item \textbf{Poids / Centrage:} Augmentent $W$ et la vitesse de rotation $V_r$ nécessaire.
				\item \textbf{Vent:} Modifie la vitesse sol initiale.
				\item \textbf{État de la piste:} C'est cette composante $W \sin \theta$ qui s'ajoute à la traînée totale et qui allonge la distance de décollage si la piste est en montée.\\
				
			\end{itemize}
		
		
	
	
	
	
	
	
	
	
	
	
	
	
	
	
	
	
	
	






